\documentclass[12pt]{article}
\usepackage[T1]{fontenc}
\usepackage[utf8]{inputenc}
\usepackage{lmodern}
\usepackage{textcomp}
\usepackage{enumitem}
\usepackage{geometry}
\usepackage{graphicx}
\usepackage{amsmath, amssymb}
\geometry{margin=1in}
\begin{document}
\begin{center}
Physics - Graduate Level Assessment 3 \
Total Marks: 40 \
Answer all questions.
\end{center}

\section*{Multiple Choice (10 marks)}
\begin{enumerate}[label=\arabic*.]
\item Among these colors, the one that has the most energy per photon is
\begin{enumerate}[label=(\alph*)]
    \item violet
    \item turquoise
    \item green
    \item blue
    \item red
\end{enumerate}
\item The red glow in the neon tube of an advertising sign is a result of
\begin{enumerate}[label=(\alph*)]
    \item fluorescence.
    \item absorption.
    \item resonance.
    \item de-excitation.
    \item refraction.
\end{enumerate}
\item A heavy rock and a light rock in free fall (zero air resistance) have the same acceleration. The heavy rock doesn't have a greater acceleration because the
\begin{enumerate}[label=(\alph*)]
    \item air resistance is always zero in free fall.
    \item volume of both rocks is the same.
    \item force due to gravity is the same on each.
    \item gravitational constant is the same for both rocks.
    \item force due to gravity is zero in free fall.
\end{enumerate}
\item When small pieces of material are assembled into a larger piece, the combined surface area
\begin{enumerate}[label=(\alph*)]
    \item greatly increases
    \item slightly increases
    \item is halved
    \item becomes zero
    \item greatly decreases
\end{enumerate}
\item A corrective lens of refractive power - 10.00 diopters corrects a myopic eye. What would be the effective power of this lens if it were 12 mm. from the eye?
\begin{enumerate}[label=(\alph*)]
    \item -11.00 diopters
    \item -9.50 diopters
    \item -10.50 diopters
    \item -9.00 diopters
    \item -7.93 diopters
\end{enumerate}
\end{enumerate}

\section*{True / False (10 marks)}
\textit{Answer True or False}
\begin{enumerate}[label=\arabic*.]
\setcounter{enumi}{5}
\item True or False: The minimum kinetic energy of an electron localized within a nuclear radius of 6 x 10^-15 m is approximately 15.9 MeV.
\item True or False: The total collisional frequency for CO 2 at 1 atm and 298 K is approximately 6.28 x $10^34$ m^-3 s^-1.
\item True or False: An electron can be effectively accelerated by a magnetic field in a straight line without any additional forces.
\item True or False: When relatively slow-moving molecules condense from the air, the temperature of the remaining air tends to decrease.
\item True or False: The complementary color of blue is cyan, which is often confused with its neighboring color on the spectrum.
\end{enumerate}

\section*{Long Form Response (20 marks)}
\textit{Answer each question with a clear and complete explanation.}
\begin{enumerate}[label=\arabic*.]
\setcounter{enumi}{10}
\item Discuss the principles of lens magnification by analyzing a scenario where a 40 cm focal length lens projects a 4 cm object onto a screen 20 m away. Calculate the size of the resulting image and explain your reasoning.
\item Discuss the energy transformations that occur during the fission of a uranium nucleus, focusing on the primary forms of energy released and their implications for nuclear physics.
\end{enumerate}

\vfill
\noindent\textit{End of Paper}
\end{document}